%%%%%%%%%%%%%%%%%%%%%%%%%%%%%%%%%%%%%%%%%%%%%%%%%%%%%%%%%%%%%%%%%%%%%%%%%%%%%%%%%
%	CHAPTER 4 - Repair Mechanism Analysis (graded focus)
%   TICKET-11 + TICKET-21
%   Owner: Theory / Writing Lead
%%%%%%%%%%%%%%%%%%%%%%%%%%%%%%%%%%%%%%%%%%%%%%%%%%%%%%%%%%%%%%%%%%%%%%%%%%%%%%%%%
\lhead{Chapter 4. \emph{Repair Mechanism Analysis}}
\chapter{Repair Mechanism Analysis}
\label{Chapter4}

\section{Sources of Infeasibility}
\label{Sec:Infeasibility}
TODO: Which operators break the permutation property; worked examples for PMX, OX, CX.

\section{Repair Effectiveness}
\label{Sec:RepairEffectiveness}
TODO: Quantitative analysis tied to experimental results.

\section{Comparison with Alternative Strategies}
\label{Sec:RepairComparison}
TODO: Repair vs.\ penalty function vs.\ feasibility-preserving operator.

\section{Impact on Search Behaviour}
\label{Sec:SearchImpact}
TODO: Discuss impact on solution quality, diversity, and convergence with reference to experimental results.
